\chapter{Four-Dimensional Superconformal Symmetry}
\label{ch:volume-iii-four-dimensional-superconformal-symmetry}

\section{The Algebraic Data}

In four Lorentzian spacetime dimensions, the conformal algebra is
\[
  \mathfrak{so}(4,2)\simeq \mathfrak{su}(2,2).
\]
Its generators are translations \(P_\mu\), Lorentz transformations
\(M_{\mu\nu}=-M_{\nu\mu}\), dilatations \(D\), and special conformal
transformations \(K_\mu\), with \(\mu,\nu=0,1,2,3\).  The
\(\mathcal N=4\) superconformal algebra has bosonic subalgebra
\[
  \mathfrak{su}(2,2)\oplus \mathfrak{su}(4)_R
\]
after quotienting by the common central element; the corresponding simple
Lie superalgebra is denoted
\[
  \mathfrak{psu}(2,2|4).
\]
The subscript \(R\) records that \(\mathfrak{su}(4)_R\) acts on the
supercharges rather than on spacetime.

Use undotted and dotted spinor indices
\[
  \alpha,\beta=1,2,
  \qquad
  \dot\alpha,\dot\beta=1,2,
\]
for the two factors of the complexified Lorentz algebra
\(\mathfrak{sl}(2,\mathbb C)\).  Use \(I,J=1,\ldots,4\) for the
fundamental representation of \(SU(4)_R\).  The Poincare supercharges are
\[
  Q^I_\alpha,
  \qquad
  \bar Q_{\dot\alpha I},
\]
and the special supercharges are
\[
  S^\alpha_I,
  \qquad
  \bar S^{\dot\alpha I}.
\]
The translations and special conformal transformations can be written as
\[
  P_{\alpha\dot\alpha}
  =
  \sigma^\mu_{\alpha\dot\alpha}P_\mu,
  \qquad
  K^{\dot\alpha\alpha}
  =
  \bar\sigma^{\mu\,\dot\alpha\alpha}K_\mu ,
\]
where \(\sigma^\mu=(1,\vec\sigma)\), \(\bar\sigma^\mu=(1,-\vec\sigma)\),
and \(\vec\sigma\) are the Pauli matrices.  The defining anticommutators
containing \(P\) and \(K\) are
\begin{equation}
  \{Q^I_\alpha,\bar Q_{\dot\beta J}\}
  =
  2\delta^I{}_J P_{\alpha\dot\beta},
  \qquad
  \{S^\alpha_I,\bar S^{\dot\beta J}\}
  =
  2\delta_I{}^J K^{\dot\beta\alpha}.
  \label{eq:n4-qq-ss-anticommutators}
\end{equation}
The mixed anticommutators \(\{Q,S\}\) are linear combinations of \(D\),
Lorentz generators, and \(SU(4)_R\) generators.  Their precise coefficients
depend on the normalization of the Cartan generators, but their structural
form is fixed:
\[
  \{Q,S\}\in
  \operatorname{span}\{D,M_{\mu\nu},R^I{}_J\}.
\]
Here \(R^I{}_J\) are traceless generators of \(\mathfrak{su}(4)_R\).

\section{Superconformal Primaries}

Radial quantization maps a local operator at the origin to a state on
\(S^3\).  The adjoint relations are
\[
  P_\mu^\dagger=K_\mu,
  \qquad
  (Q^I_\alpha)^\dagger=S^\alpha_I,
  \qquad
  (\bar Q_{\dot\alpha I})^\dagger=\bar S^{\dot\alpha I}.
\]
A superconformal primary is a state \(\ket{\mathcal O}\) satisfying
\begin{equation}
  K_\mu\ket{\mathcal O}=0,
  \qquad
  S^\alpha_I\ket{\mathcal O}=0,
  \qquad
  \bar S^{\dot\alpha I}\ket{\mathcal O}=0.
  \label{eq:superconformal-primary-definition}
\end{equation}
It is also chosen to be a highest-weight state for
\[
  SU(2)_{\rm left}\times SU(2)_{\rm right}\times SU(4)_R .
\]
Its quantum numbers are
\[
  \Delta,\qquad
  (j,\bar j),\qquad
  [r_1,r_2,r_3],
\]
where \(\Delta\) is the eigenvalue of \(D\), \(j,\bar j\) are the two
Lorentz spins, and \([r_1,r_2,r_3]\) are Dynkin labels of \(SU(4)_R\).
The integers \(r_i\) are nonnegative.

Descendants are obtained by applying \(P_\mu\), \(Q^I_\alpha\), and
\(\bar Q_{\dot\alpha I}\).  Because the supercharges are odd generators,
their action creates a finite fermionic envelope over the ordinary conformal
descendants.

\begin{figure}[H]
  \centering
  \resizebox{0.92\textwidth}{!}{%
  \begin{tikzpicture}[x=1cm,y=1cm,every node/.style={font=\scriptsize}]
    \tikzset{
      arrow/.style={-{Latex[length=2mm]},line width=0.85pt},
      insertion/.style={circle,fill=violet!80!black,inner sep=1.8pt},
      statebox/.style={rounded corners=4pt,draw,line width=0.8pt,fill=blue!4,
        minimum width=2.65cm,minimum height=1.00cm}
    }
    \node[statebox] (primary) at (-3.9,0) {primary \(\ket{\mathcal O}\)};
    \node[statebox] (qdesc) at (-0.3,1.05) {\(Q\)-descendants};
    \node[statebox] (pdesc) at (-0.3,-1.05) {\(P\)-descendants};
    \node[statebox] (short) at (3.5,0) {short multiplet};
    \draw[arrow] (primary)--(qdesc) node[midway,above left] {\(Q,\bar Q\)};
    \draw[arrow] (primary)--(pdesc) node[midway,below left] {\(P_\mu\)};
    \draw[arrow,green!45!black] (qdesc)--(short);
    \node[green!45!black,fill=white,inner sep=1pt] at (2.38,0.92) {null states removed};
    \draw[arrow] (pdesc)--(short);
    \node[align=center,blue!65!black] at (0.0,0)
      {\(\{Q,S\}\) gives\\norm constraints};
    \node[align=center,violet!80!black] at (3.5,-1.30)
      {shortening is an algebraic\\property of the representation};
  \end{tikzpicture}
  }
  \caption{A superconformal multiplet is generated from a superconformal
  primary by translations and Poincare supercharges.  Shortening occurs when
  some supercharge descendants have zero norm and are quotiented out.}
  \label{fig:superconformal-primary-descendants}
\end{figure}

\section{Positivity and Shortening}

Let \(q\) be any complex linear combination of Poincare supercharges and let
\(q^\dagger\) be its radial-quantization adjoint.  For a primary state,
\[
  \|q\ket{\mathcal O}\|^2
  =
  \bra{\mathcal O}q^\dagger q\ket{\mathcal O}
  =
  \bra{\mathcal O}\{q^\dagger,q\}\ket{\mathcal O}.
\]
The last expression is an eigenvalue of a Hermitian linear combination of
\(D\), Lorentz Cartan generators, and \(SU(4)_R\) Cartan generators.  Unitarity
therefore gives inequalities of the form
\begin{equation}
  \Delta\ge L_q(j,\bar j;r_1,r_2,r_3),
  \label{eq:superconformal-positivity-bound}
\end{equation}
where \(L_q\) is the linear functional determined by the weight of the
supercharge \(q\).  When equality holds for some nonzero \(q\), the state
\(q\ket{\mathcal O}\) has zero norm.  The corresponding descendant is removed
from the irreducible Hilbert-space representation.  This is a shortened
multiplet.

A particularly important family consists of scalar primaries with Lorentz
spins \((0,0)\) and \(SU(4)_R\) Dynkin labels \([0,p,0]\).  In the standard
normalization of \(SU(4)_R\), positivity gives
\begin{equation}
  \Delta\ge p.
  \label{eq:half-bps-bound}
\end{equation}
Saturation is the half-BPS condition.  It is equivalent to annihilation by
half of the Poincare supercharges and their adjoints in the conjugate
directions.  A half-BPS primary therefore generates a strictly smaller
superconformal multiplet than a generic primary with the same \(SU(4)_R\)
highest weight.

\section{The R-Symmetry Weights}

The group \(SU(4)_R\) has rank three.  An irreducible finite-dimensional
representation is labeled by a highest weight with Dynkin labels
\[
  [r_1,r_2,r_3],
  \qquad
  r_i\in\mathbb Z_{\ge0}.
\]
The fundamental representations used below are
\[
  \mathbf4=[1,0,0],
  \qquad
  \overline{\mathbf4}=[0,0,1],
  \qquad
  \mathbf6=[0,1,0],
\]
and the adjoint representation is
\[
  \mathbf{15}=[1,0,1].
\]
The six-dimensional representation \(\mathbf6\) is the vector
representation of \(SO(6)\), using the isomorphism
\[
  SU(4)_R\simeq \operatorname{Spin}(6).
\]

\begin{figure}[H]
  \centering
  \resizebox{0.88\textwidth}{!}{%
  \begin{tikzpicture}[x=1cm,y=1cm,every node/.style={font=\scriptsize}]
    \tikzset{arrow/.style={-{Latex[length=2mm]},line width=0.85pt}}
    \node[font=\small] at (0,1.85) {\(SU(4)_R\) Dynkin labels};
    \foreach \x/\lab in {-2.4/{r_1},0/{r_2},2.4/{r_3}} {
      \draw[thick] (\x,0) circle (0.42);
      \node at (\x,0) {\(\lab\)};
    }
    \draw[thick] (-1.98,0)--(-0.42,0);
    \draw[thick] (0.42,0)--(1.98,0);
    \node[align=center,blue!65!black] at (-2.4,-1.15)
      {\(\mathbf4\)\\\([1,0,0]\)};
    \node[align=center,green!45!black] at (0,-1.15)
      {\(\mathbf6\)\\\([0,1,0]\)};
    \node[align=center,violet!80!black] at (2.4,-1.15)
      {\(\overline{\mathbf4}\)\\\([0,0,1]\)};
    \draw[arrow,orange!85!black] (3.25,0.35)--(4.35,0.35);
    \node[right,align=left] at (4.42,0.35)
      {half-BPS scalar family\\\([0,p,0]\), \(\Delta=p\)};
  \end{tikzpicture}
  }
  \caption{The \(SU(4)_R\) representation of a superconformal primary is
  recorded by three Dynkin labels.  The half-BPS scalar family has labels
  \([0,p,0]\) and protected dimension \(\Delta=p\).}
  \label{fig:su4-dynkin-bps-family}
\end{figure}

The \(p=2\) member of the half-BPS family is the stress-tensor multiplet
primary.  It will reappear as the protected scalar operator in
\(\mathcal N=4\) super-Yang--Mills.
